% Directly inputtable article fragment.
% Required preamble packages: amsmath, amssymb, amsthm.
% Author: ChatGPT, 2 September 2026.

\section{All-index Pell polynomials and complete Hensel steering}
\label{sec:pell-all-index-formalization}

\noindent\textbf{Author: ChatGPT.}\quad
This section closes the algebraic formalization gaps left in
Section~\ref{sec:pell-polynomial-hensel}.  It does not prove the fixed
coefficient-two squarefull exclusion and does not prove or disprove the
$abc$ conjecture.  Its purpose is to put the all-index identities, the full
polynomial Hensel law, the all-support equivalence, finite simultaneous
steering, and the complete moving counterexample under one exact interface.

Let $F_n,L_n\in\mathbb Z[T]$ be defined by
\[
 F_0=0,\quad F_1=1,\quad F_{n+2}=TF_{n+1}+F_n,
\]
\[
 L_0=2,\quad L_1=T,\quad L_{n+2}=TL_{n+1}+L_n.
\]

\begin{theorem}[All-index polynomial identities]
\label{thm:pell-all-index-polynomial}
For every $n\geq0$,
\begin{align}
 L_n^2-(T^2+4)F_n^2&=4(-1)^n,\label{eq:pell-all-index-norm}\\
 L_n'&=nF_n,\label{eq:pell-all-index-L-derivative}\\
 (T^2+4)F_n'&=nL_n-TF_n.\label{eq:pell-all-index-F-derivative}
\end{align}
Evaluation at any $t\in\mathbb Z$ preserves the two recurrences and all
three identities.
\end{theorem}

\begin{proof}
Two-step induction gives
\begin{equation}\label{eq:pell-companion-relations}
 L_n=2F_{n+1}-TF_n,
 \qquad L_{n+1}=TF_{n+1}+2F_n.
\end{equation}
Ordinary induction in the $F$ recurrence gives Cassini's identity
\[
 F_{n+1}^2-TF_{n+1}F_n-F_n^2=(-1)^n.
\]
Substitution of the first equality in
\eqref{eq:pell-companion-relations} proves
\eqref{eq:pell-all-index-norm}.  Differentiating the $L$ recurrence and
using the second equality in \eqref{eq:pell-companion-relations} proves
\eqref{eq:pell-all-index-L-derivative} by two-step induction.  For the
remaining formula, differentiate the $F$ recurrence and use its two
preceding instances together with
\[
 (T^2+4)F_{n+1}=TL_{n+1}+2L_n,
\]
which follows from \eqref{eq:pell-companion-relations}.  Polynomial
evaluation is a ring homomorphism, so it transports every displayed equality.
\end{proof}

\begin{theorem}[Polynomial Taylor--Hensel law]
\label{thm:pell-full-polynomial-hensel}
Let $f\in\mathbb Z[T]$, let $p,t,c,h\in\mathbb Z$, and let $e\geq1$.
For every fixed $z\in\mathbb Z$ there is an integer
$G=G(f,t,z)$ such that
\[
 f(t+z)=f(t)+zf'(t)+z^2G.
\]
In particular,
\[
 f(t+p^eh)\equiv f(t)+p^ehf'(t)\pmod {p^{e+1}}.
\]
If $p\neq0$ and $f(t)=p^ec$, then
\begin{equation}\label{eq:pell-hensel-divisibility-iff}
 p^{e+1}\mid f(t+p^eh)
 \quad\Longleftrightarrow\quad
 p\mid c+hf'(t).
\end{equation}
If moreover $\gcd(f'(t),p)=1$, a digit satisfying
\eqref{eq:pell-hensel-divisibility-iff} exists and is
unique modulo $p$.
\end{theorem}

\begin{proof}
The integral Taylor expansion gives the first formula.  Substitution of
$z=p^eh$ and the inequality $e+1\leq2e$ give the congruence.  After writing
$f(t)=p^ec$, cancellation of the nonzero factor $p^e$ gives
\eqref{eq:pell-hensel-divisibility-iff}.
B\'ezout's identity supplies an admissible digit, and subtracting two
admissible affine congruences proves uniqueness.
\end{proof}

Call $N\in\mathbb N$ squarefull when every prime $p\mid N$ satisfies
$p^2\mid N$.  For a polynomial $f$, an integer $t$, and a support prime
$p$, define zero first displacement to mean
\begin{equation}\label{eq:pell-zero-first-displacement}
 p\mid f(t),\qquad \gcd(f'(t),p)=1,
 \qquad p^2\mid f(t).
\end{equation}
By Theorem~\ref{thm:pell-full-polynomial-hensel}, the final condition is
equivalent to the unique next Hensel digit being zero.

\begin{proposition}[All-support two-channel equivalence]
\label{prop:pell-all-support-equivalence}
Let $A,B\in\mathbb N$ be coprime, let
$f_A(t)=u_AA$ and $f_B(t)=u_BB$, and suppose that each scale $u_A,u_B$ is
coprime to every prime in the corresponding support.  Assume also that all
support roots of $f_A,f_B$ at $t$ are simple.  Then $AB$ is squarefull if
and only if \eqref{eq:pell-zero-first-displacement} holds at every support
prime in both channels.

In particular, this applies formally to
\[
 f_A=L_\ell,\quad u_A=2,\quad
 f_B=F_\ell,\quad u_B=1,\quad t=2,
\]
once the odd-support scale condition and both all-support transversality
conditions have been proved.
\end{proposition}

\begin{proof}
For coprime $A,B$, Euclid's lemma gives
\[
 AB\text{ squarefull}
 \quad\Longleftrightarrow\quad
 A\text{ and }B\text{ squarefull}.
\]
In each channel, a support prime square divides $u_AN$ if and only if it
divides $N$, since it is coprime to the scale.  The asserted equivalence now
follows directly from \eqref{eq:pell-zero-first-displacement}.
\end{proof}

\begin{theorem}[Finite simultaneous Hensel steering]
\label{thm:pell-finite-indexed-steering}
Let $f_1,\ldots,f_m\in\mathbb Z[T]$, let $p_1,\ldots,p_m$ be distinct
primes, and suppose
\[
 p_i\mid f_i(t_0),\qquad p_i\nmid f_i'(t_0)
 \quad(1\leq i\leq m).
\]
There is a unique residue class
\[
 t\pmod{\prod_{i=1}^m p_i^2}
\]
such that
\[
 t\equiv t_0\pmod {p_i},\qquad p_i^2\mid f_i(t)
 \quad(1\leq i\leq m).
\]
\end{theorem}

\begin{proof}
Theorem~\ref{thm:pell-full-polynomial-hensel} gives one and only one local
class modulo $p_i^2$.  The square moduli are pairwise coprime, so the integer
Chinese remainder theorem gives a joint class and uniqueness modulo their
product.  Polynomial evaluation respects congruence.  Conversely, any
integer satisfying the displayed base and divisibility conditions lies in
each unique local class, hence in the joint CRT class.
\end{proof}

At index three the recurrence gives
\[
 F_3(T)=T^2+1,\qquad L_3(T)=T^3+3T.
\]
The selected roots at $t_0=2$ are
\[
 7\mid L_3(2),\quad 7\nmid L_3'(2),\qquad
 5\mid F_3(2),\quad 5\nmid F_3'(2).
\]
Their lifted classes are $t\equiv37\pmod {49}$ and
$t\equiv7\pmod {25}$; the explicit joint representative used below is $282$.

\begin{proposition}[Complete moving witness]
\label{prop:pell-complete-moving-witness}
Put
\[
 t=282,\quad A=11213307,\quad B=79525,\quad D=19882.
\]
Then
\[
 2A=L_3(t),\qquad B=F_3(t),\qquad 4D=t^2+4,
\]
\[
 A=3^2\,7^2\,47\,541,\qquad
 B=5^2\,3181,\qquad D=2\,9941,
\]
all displayed factors are prime, $D$ is positive and squarefree, and
\[
 A^2-DB^2=-1.
\]
The selected factors satisfy $7^2\mid A$ and $5^2\mid B$, while
\[
 47^2\nmid A,\qquad541^2\nmid A,
 \qquad3181^2\nmid B.
\]
Thus this tuple satisfies every premise of the moving-parameter exclusion
and refutes that exact strengthening.  It is not a squarefull two-channel
packet and has $D\neq2$.
\end{proposition}

\begin{proof}
All evaluations and factorizations are direct integer identities.  Trial
division proves the primality assertions; in particular $9941$ is prime,
so $D=2\cdot9941$ is squarefree.  The global norm is the specialization of
\eqref{eq:pell-all-index-norm}, divided by four.  The three failed square
divisibilities certify the stated non-squarefull scope.
\end{proof}

The companion Lean module checks each result above.  Its structure
\texttt{HGlobalMoveWitness} contains the prime-index,
base-root, simple-derivative, lifted-residue, repeated-divisibility,
factorization, primality, positivity, squarefree-coefficient, global-norm,
and non-squarefull fields in one object.  The associated axiom audit reports
only the standard Mathlib principles \texttt{propext},
\texttt{Classical.choice}, and \texttt{Quot.sound}.

The arithmetic route remains open at a precise point: one must discharge the
all-support transversality premises for the actual fixed Pell packet and then
exclude simultaneous zero displacements at $T=2$.  No counterexample meeting
those full fixed-parameter premises has been found, so the route is retained.
